\documentclass[11pt,a4paper]{article}
\usepackage[margin=2.2cm]{geometry}
\usepackage{amsmath,amssymb,xcolor,enumitem,booktabs,graphicx}
\usepackage[normalem]{ulem}
\setlength{\parskip}{0.6em}\setlength{\parindent}{0pt}
\definecolor{newc}{HTML}{1f6f43}
\newcommand{\new}[1]{\textcolor{newc}{#1}}
\newcommand{\old}[1]{\textcolor{gray}{\sout{#1}}}
\newcommand{\note}[1]{{\small\itshape [#1]}}
\newcommand{\lab}[1]{\textsf{\footnotesize #1}}
\renewcommand{\ref}[1]{\lab{#1}}
\renewcommand{\eqref}[1]{(\lab{#1})}
\newcommand{\nref}[1]{the (\lab{#1})}
\newcommand{\citep}[1]{[\lab{#1}]}
\newcommand{\citet}[1]{\lab{#1}}
\title{Proposed change, ``When moves still clash''}
\date{24 November 2026, version 1}
\begin{document}\sloppy
\maketitle

\section*{How to read this}
This answers what is left of \texttt{method\_resolution\_comments.md} for
\lab{sec:method:resolution} in \texttt{4.Implementation.tex}. A3 (the lock
told as agents choosing to wait), C1, C2 and the dock draw of C3 were done in
[A107] and [A109]. This change does D1, D2 (keep the list, cut the prose, as
agreed on 24 November), D3, C3 and C4, and adds parking to the legend of
Figure~15. The rule itself, the code and the lock story do not change.

Black text stays, \new{green} is new and \old{struck} text goes.
\note{Notes in brackets} say why. Nothing below is in the thesis yet.

\section*{Part 1. What was checked}
\begin{itemize}[nosep]
  \item \texttt{resolution/conflict\_resolution.py}: an agent whose chain of
    waiting ``runs into a cycle'' is never ready and waits (D3).
  \item The Problem Formulation (Who decides) already says the priority
    order is drawn once per episode.
  \item \texttt{styles/eqnames.tex}: \lab{eq:waiting} is named ``mean
    blocked time'', so ``the blocked time of \texttt{\textbackslash
    nref\{eq:waiting\}}'' prints ``the blocked time of the mean blocked time
    (31)'' (C4).
  \item Figure~15 already follows the viewer: orange to a pick, red to a
    workstation. The green diamonds of (a) are parking spots and are missing
    from the legend.
\end{itemize}

\section*{Part 2. The section as it would read}

\old{The masks make each move possible on its own, but two agents that
choose separately can still pick the same vertex or try to swap.}
\new{The action mask makes each move possible on its own. Inside one
decision of the centralised and section-based controllers the claim mask
also keeps the agents apart. Agents that choose separately can still pick the
same vertex or try to swap. These are decentralised agents, which choose at
the same time, and agents of two zones at a border.} A penalty in the reward
would not be enough on its own, because a learned policy draws its actions
at random and can still draw a forbidden move. A fixed rule removes these
clashes. \old{It goes through the agents one at a time, accepts a move that
does not clash with the moves already accepted, and turns any other move into
a wait.} An agent that wants the vertex of another agent is handled after
that agent, so it can follow an agent that leaves, as in
\citet{vanKnippenberg2021}. \old{Among the agents that are ready, the
priority order decides, and it is drawn once from the episode seed and kept
for the whole episode.} Formally, given the proposed joint action
$(u_1(t),\dots,u_m(t))$, agent $a_i$ \emph{waits on} $a_j$ when $u_i(t)$
moves onto $\ell_j(t)$, and
\begin{enumerate}[nosep]
  \item an agent is ready when it waits on no agent, or on an agent already
    handled. Of the ready agents, the first in the priority order is
    handled next.
  \item the vertex it asks for is accepted if it breaks neither the vertex
    rule nor the swap rule against the moves already accepted.
  \item otherwise its action becomes $\mathrm{wait}$ (always legal by
    \nref{eq:actions}) and its override flag for $t$ is set.
  \item agents never ready wait, with their override flag set.
\end{enumerate}
\note{D1 and D2. The list is the rule. The prose kept only what the list does
not say, why a fixed rule is needed and why the agent in front goes first.
The priority order is drawn once per episode in the Problem Formulation.}

An agent is never ready when the agents it waits on form a circle, each
wanting the next one's vertex\new{, or when its chain of waiting runs into
such a circle}. A swap is the shortest such circle. These agents wait, since
none of them can move first. \note{D3.}

\emph{[The lock, ``No rule is added to break the lock'', Figure~15 and its
story are unchanged.]}

\textbf{Figure 15 legend.} shelf, workstation, dock, \new{parking,} agent to
a pick, agent to a workstation. \note{Script
\texttt{scripts/images/workstation\_lock.py}, one entry added with the
viewer's green diamond.}

The rule always ends with no collision. \emph{[Unchanged.]}

The rule is a fixed function of the state, which includes the priority order
(Section~\ref{sec:pf:game}), with no added randomness. So \old{$P$}
\new{the warehouse} is random only through the arrival of orders and the dock
drawn at random for each shipment and put-away. \old{It also makes paired
seeds a fair comparison, because the same seed and the same moves give the
same decisions.} \new{So two runs with the same seed and the same moves make
the same decisions, which is what a comparison on the same seeds needs.}
\note{C3. $P$ was last met chapters ago, and ``paired seeds'' is defined only
in the Experiments.}

The override flag is logged (Section~\ref{sec:impl:logging}) and used in
two places. It sets the override term of \nref{eq:reward}, and together
with a chosen wait it \old{sets the blocked time of \nref{eq:waiting}}
\new{counts towards the mean blocked time~\eqref{eq:waiting}}. \note{C4.}

\section*{Questions}
None. The choices were made in chat on 24 November.

\end{document}
